% ============================================================================
% Единое Естествознание — Книга II: Лабораторный практикум ДКМ
% (book/sources/volumes/vol-laboratory-manual.tex)
% ============================================================================

\chapter*{Введение. Лабораторный арбитраж в едином Естествознании и связь четырёх книг}
\addcontentsline{toc}{chapter}{Введение. Лабораторный арбитраж в едином Естествознании и связь четырёх книг}

В рамках дискретной квантовой механики (ДКМ) на двусортной алгебраической системе носителей А.\,И.~Мальцева $\langle\mathcal{M},\mathcal{T}\rangle$ ($\mathcal{M} = \mathbb{Z}_{512}$ на оболочке кубооктаэдра $O_2(3)$ при $1 + 2^3 = 3^2$) физика, химия, молекулярная биология, космология и физическая реализация вычислений объединяются в \textbf{единое Естествознание}, изложенное в четырёх взаимно дополняющих книгах без дублирования материала:
\begin{enumerate}[label=\textbf{Книга~\Roman*.},leftmargin=*,noitemsep]
\item \textbf{Теория (\texttt{formal-foundations-monograph.pdf}):} строгий дедуктивный вывод законов ДКМ и непрерывной квантовой механики (НКМ), уравнений Максвелла, постоянной тонкой структуры $\alpha_{\mathrm{fs}}^{-1}(3) = 137{,}035\,999\,187\dots$, атомных и ядерных оболочек $O_2(3)$, молекулярной биофизики и теоретико-модельного предела $\lim_{n\to\infty} M_n = \mathcal{M}$ (\texttt{ModelTheoreticLimit.lean}, \texttt{UnifiedOperators.lean}).
\item \textbf{Лаборатория (\texttt{dqm-laboratory-manual.pdf}, настоящий том):} пошаговые лабораторные протоколы чистой комнаты, контрольные A/B-эксперименты, количественные сигнатуры фальсификации (\texttt{M1}--\texttt{M4}), упреждающая диагностическая матрица (\texttt{F1}--\texttt{F6}) и журнал критических экспериментов, выполняющие роль внешнего семантического арбитра по теореме Ловера--Гёделя--Тарского (\texttt{thm17}).
\item \textbf{Инженерия (\texttt{dqm-engineering-handbook.pdf}):} теоретическая электротехника фазовой эстафеты ($V = \mathcal{C}_3(\Delta\varphi),\ I = \mathcal{S}_3(\Delta\varphi)$), беспульсационные квадратурные выпрямители, физическая реализация вычислений без тепла Ландауэра в коридоре Гейзенберга $|\Delta\varphi|\le 40 < 41$, кремниевая архитектура \texttt{Ring9} / \texttt{Bus36} / \texttt{VerletCell72} и 12-лучевые резонаторы~$O_2(3)$.
\item \textbf{Научно-популярный путеводитель (\texttt{natural-sciences-course.pdf}, \emph{Powered by Dr.\,STONE}):} разговорный курс и дорожная карта всех семи ступеней крафта Вселенной.
\end{enumerate}
Историческая хроника объединения восьми отдельных дисциплин («Фильм о фильме / Как это создавалось») сохранена в каталоге \texttt{book/out/pdf/making-of/}.

% ============================================================================
\chapter{Бесфононная кремниевая фотовольтаика (\texorpdfstring{$441\text{ нм} \to 2\times 882\text{ нм} < 1107\text{ нм}$}{441 nm -> 2 x 882 nm < 1107 nm}): протокол чистой комнаты и матрица диагностики}
\label{ch:lab-solar-photovoltaics}
% ============================================================================

\section{Физическая причина стагнации однопороговой солнечной энергетики: несогласованная фононная термализация}

В стандартной кремниевой фотовольтаике ($\mathrm{Si}$, $Z = 14$) полупроводниковый переход работает как однопороговый фильтр с краем поглощения
\begin{equation}
  \label{eq:lab-si-bandgap}
  \lambda_{\mathrm{Si}}(3) \;=\; d^d \cdot \Delta\varphi_{\mathrm{disc}}(3) \;=\; 27 \times 41 \;=\; 1\,107\text{ нм} \qquad (E_g \approx 1{,}12\text{ эВ}).
\end{equation}
Фундаментальный предел Шокли--Квайссера ($\sim 33{,}7\%$ для однопереходного элемента) обусловлен тем, что видимый и сине-фиолетовый спектр солнечного излучения ($\lambda \in [400, 720]\text{ нм}$) поглощается кремнием в режиме жёсткого рассогласования импедансов:
\begin{enumerate}[leftmargin=*,noitemsep]
\item Синий фотон в главном планетарном окне $\lambda_{\mathrm{blue}} = \Delta\varphi_{\mathrm{seam}}(3)^2 = 21^2 = \mathbf{441\text{ нм}}$ ($E \approx 2{,}81\text{ эВ}$) обладает энергией, более чем вдвое превышающей ширину запрещённой зоны кремния $E_g \approx 1{,}12\text{ эВ}$.
\item В обычном кремнии такой фотон рождает всего одну электрон-дырочную пару, а избыток энергии ($\sim 1{,}69\text{ эВ}$, или $60\%$ энергии фотона) сбрасывается в кристаллическую решётку в виде надпороговых тепловых фононов ($\hat{n}_E > 0$).
\item Выделившееся паразитное тепло разогревает солнечную панель под инсоляцией до $338\text{--}348\text{ К}$, снижая напряжение холостого хода~$V_{\mathrm{oc}}$ на $\sim 0{,}4\%$ на каждый кельвин нагрева.
\end{enumerate}

\section{Бесфононное деление синего кванта (Singlet Fission: \texorpdfstring{$441\text{ нм} \to 2 \times 882\text{ нм} < 1107\text{ нм}$}{441 nm -> 2 x 882 nm < 1107 nm})}

В процессе деления синглетного экситона (\emph{Singlet Fission}, SF) синглетный квант $S_1$ делится на пару спин-связанных триплетных экситонов $T_1 + T_1$ с удвоенной длиной волны:
\begin{equation}
  \label{eq:lab-sf-882}
  \lambda_{\mathrm{triplet}}(3) \;=\; \left(\frac{c_0}{c}\right)^2 \lambda_{\mathrm{blue}}(3) \;=\; 2 \times 441 \;=\; \mathbf{882\text{ нм}} \;=\; |\square(3)|\cdot |B_3(3)| \;=\; 6 \times 147,
\end{equation}
\begin{equation}
  \label{eq:lab-sf-margin}
  \Delta\lambda_{\mathrm{Si - SF}}(3) \;=\; \lambda_{\mathrm{Si}}(3) - \lambda_{\mathrm{triplet}}(3) \;=\; 1\,107 - 882 \;=\; \mathbf{225\text{ нм}} \;=\; 15^2.
\end{equation}
Оба триплетных экситона ($882\text{ нм}$, $1{,}406\text{ эВ}$) лежат строго внутри окна поглощения кремния с запасом $15^2 = 225\text{ нм}$.

\section{Согласование оболочечного и фазового импеданса на границе «Тетрацен -- \texorpdfstring{$\mathrm{SiO_x}$}{SiOx} -- Кремний»}

Полное доказательство теоремы \texttt{checkZeroPhononSingletFissionSiliconPhotovoltaicBoostOp} (\texttt{lean/MtCarrier/UnifiedOperators.lean}) приведено в Книге~I; здесь зафиксируем четыре рабочих инварианта для инженера чистой комнаты:
\begin{enumerate}[label=\textbf{\arabic*.},noitemsep]
\item \textbf{Оболочечный дефицит тетрацена на голом кремнии:}
  $Z(\mathrm{C_{18}H_{12}}) = 18\times 6 + 12\times 1 = \mathbf{120} = Q_\square(3) - |{\triangle}(3)| = 128 - 8$. На прямом контакте с голым $\mathrm{Si}(100)$ ($Z=14$) оболочка не достаёт до седла $Q_\square = 128$ на $8$ единиц, а сдвиг синглет-триплетной щели равен $|N(3)|^2 = 12^2 = \mathbf{144\text{ мэВ}}$ (\texttt{arXiv:2310.15736}).
\item \textbf{Кислородно-силанольный замок седла $120 + 8 = 128$ и барьер Паули $46^2 > 2048$:}
  кислородный мост $Z(\mathrm{O}) = 8$ в составе $\mathrm{SiO_4}$ ($Z = 14 + 32 = 46$) или водородно-пассивированного $[\mathrm{AlO_4H}]$ ($Z = 13 + 32 + 1 = 46$) выводит тетрацен точно на седло ротатора $120 + 8 = 128 = Q_\square(3)$ ($\mathcal{C}_3(128) = 0,\ \mathcal{S}_3(128) = +128$) и запирает неразделённый синглет $S_1$ порогом Паули:
  \begin{equation}
    \label{eq:lab-pauli-passivation}
    Z(\mathrm{AlO_4})^2 = 45^2 = 2\,025 \;<\; 2\,048 \;<\; 2\,116 = 46^2 = Z(\mathrm{SiO_4})^2 = Z([\mathrm{AlO_4H}])^2.
  \end{equation}
\item \textbf{Подавление Оже-рекомбинации и просветление Френеля $288/289$:}
  разлёт триплетов по ортогональным лучам $O_2(3)$ ($\Delta\varphi = 128 \implies \mathcal{C}_3(128) = 0$) блокирует биэкситонную Оже-аннигиляцию (\texttt{arXiv:2606.31359}), шаг переноса $\Delta\varphi_{\mathrm{seam}} = 21 < 41$ лежит внутри бесфононного коридора Гейзенберга ($\hat{n}_E(21) = 0$), а оптическое согласование $n_1/n_2 = 9/8$ даёт $T_{\mathrm{Fresnel}} = 288/289 \approx 99{,}654\%$.
\item \textbf{Замыкание каскада через вронскиан Мейсона--Стотерса--Безу на Замке~IV ($3\cdot 1^3 + 1\cdot 5^3 = 2\cdot 4^3 = 128$):}
  при $(X,Y,Z) = (1,5,4)$ имеем $5^3 - 5 = 120 = Z(\mathrm{C_{18}H_{12}})$, $\operatorname{Tr}\mathcal{M}_3^{\mathrm{adj}} = 4^2 + 128 = 144$, и $-\det\mathcal{W}_{\mathrm{MS}} = 1^4 + 5^4 + 4^4 = \mathbf{882\text{ нм}} = \lambda_{\mathrm{triplet}}(3)$.
\end{enumerate}

\section{Пошаговый протокол изготовления в чистой комнате и контрольный A/B-эксперимент}

\begin{table}[htbp]
\centering
\caption{Четырёхслойный бесфононный каскад солнечной ячейки с согласованным импедансом ($d=3$).}
\label{tab:lab-solar-cascade}
\small
\begin{tabular}{@{}llll@{}}
\toprule
\textbf{Слой каскада} & \textbf{Материал и геометрия} & \textbf{Инвариант ДКМ} & \textbf{Физическая функция} \\
\midrule
1. Входная оптика & Кварцевый микрорельеф $O_2(3)$ & $n_1/n_2 = 9/8$, $T = 288/289$ & Просветление $99{,}65\%$ и захват 12 лучей \\
2. Делитель кванта & Кристаллический тетрацен $\mathrm{C_{18}H_{12}}$ & $Z = 120$, $441\text{ нм} \to 2\times 882\text{ нм}$ & Деление синего фотона $S_1 \to T_1 + T_1$ ($\hat{n}_E=0$) \\
3. Седловой мост & Оксид/силанол $\mathrm{SiO_4} / [\mathrm{AlO_4H}]$ & $120 + 8 = 128 = Q_\square$, $Z = 46$ & Барьер Паули $46^2 > 2048$ для $S_1$, окно для $T_1$ \\
4. Базовый переход & Кремний IBC-$\mathrm{Si}$ (тыльные контакты) & $Z = 14$, $\lambda_{\mathrm{Si}} = 27\times 41 = 1107\text{ нм}$ & Приём $2\times 882\text{ нм}$ с запасом $15^2 = 225\text{ нм}$ \\
\bottomrule
\end{tabular}
\end{table}

\begin{enumerate}[label=\textbf{Шаг~\arabic*.},leftmargin=*]
\item \textbf{Базовая подложка (Слой~4):} пластина кристаллического кремния $n$-типа $\mathrm{Si}(100)$ со встречно-штыревыми тыльными контактами (IBC-$\mathrm{Si}$, $Z(\mathrm{Si})=14$, $\lambda_{\mathrm{Si}} = 1107\text{ нм}$, удельное сопротивление $1\text{--}5\ \Omega\cdot\text{см}$), очистка RCA + HF-dip до атомарно гладкой грани $\mathrm{Si}(100)$.
\item \textbf{Седловой пассивирующий мост и контрольный A/B-тест порога Паули (Слой~3):} выращивание ультратонкого ($0{,}8\text{--}1{,}2\text{ нм}$, $2\text{--}3$ координационных монослоя) туннельного диэлектрического моста на передней грани $\mathrm{Si}(100)$ в двух режимах:
  \begin{itemize}[noitemsep]
  \item \emph{Контрольный образец~A ($Z = 45$, затвор Паули открыт для утечки):} дегидрированный сухой оксид $\mathrm{AlO_4}$ ($Z = 13 + 32 = 45$, вакуумный отжиг $> 670\text{ К}$, $45^2 = 2025 < 2048 = K_P(3)^2/2$), допускающий сброс неразделённого синглета $S_1$;
  \item \emph{Целевой образец~B ($Z = 46$, седловой замок Паули):} тетраэдрический силанол/кремнезём $\mathrm{SiO_4}$ ($Z = 14 + 32 = 46$) либо водородно-пассивированный ALD $[\mathrm{AlO_4H}]$ ($Z = 13 + 32 + 1 = 46$, отжиг в форминг-газе $\mathrm{H_2}/\mathrm{N_2}$ при $520\text{--}570\text{ К}$, $46^2 = 2116 > 2048$), запирающий неразделённый синглет $S_1$ и выводящий тетрацен ($Z=120$) через терминальный атом кислорода ($Z=8$) точно на седло четверти кольца $120 + 8 = 128 = Q_\square(3)$.
  \end{itemize}
\item \textbf{Нанесение слоя деления синглетов (Слой~2):} термическое напыление в вакууме ($\le 10^{-6}\text{ Па}$) сублимированного тетрацена $\mathrm{C_{18}H_{12}}$ ($Z = 120$, толщина $30\text{--}70\text{ нм}$, скорость $0{,}2\text{--}0{,}5\text{ \AA/с}$ при температуре подложки $295\text{--}300\text{ К}$) на седловой мост с ориентацией осей вдоль ортогональных лучей $O_2(3)$ ($\Delta\varphi = 128 \implies \mathcal{C}_3(128) = 0$).
\item \textbf{Входная просветляющая инкапсуляция (Слой~1):} кварцевое покрытие ($\mathrm{SiO_2}$, $Z = 30$) с 12-лучевым микрорельефом кубооктаэдра $O_2(3)$ и отношением показателей преломления $n_1/n_2 = 9/8$ ($T_{\mathrm{Fresnel}} = 288/289 \approx 99{,}654\%$), наносимое без выноса образца на воздух.
\end{enumerate}

\section{Четыре сигнатуры фальсификации (\texttt{M1}--\texttt{M4}) и упреждающая матрица коррекции (\texttt{F1}--\texttt{F6})}

\begin{enumerate}[label=\textbf{(M\arabic*)},leftmargin=*]
\item \textbf{Спектроскопия границы раздела (UPS / IPES / TR-PL):} обменный сдвиг синглет-триплетной щели $\operatorname{Tr}\mathcal{M}_3^{\mathrm{adj}} = 4^2 + 128 = \mathbf{144\text{ мэВ}}$ и пик фосфоресценции/переноса триплетного экситона на $-\det\mathcal{W}_{\mathrm{MS}} = 1^4 + 5^4 + 4^4 = \mathbf{882\text{ нм}}$ ($1{,}406\text{ эВ}$).
\item \textbf{Внешний квантовый выход ($\mathrm{EQE}$) и магнито-фототок ($\mathrm{MFE}$):} подъём $\mathrm{EQE}(441\text{ нм}) > 100\%$ на синем пике тетрацена с характерной отрицательной модуляцией фототока в магнитном поле $B \in [0,\ 0{,}4]\text{ Тл}$ в Целевом образце~B ($Z=46$) при отсутствии подъёма в Контрольном образце~A ($Z=45$).
\item \textbf{Безожевая линейность от интенсивности:} линейный рост триплетного вклада в фототок в диапазоне $0{,}1\text{--}5\text{ Sun}$ без квадратичного Оже-завала ($\mathcal{C}_3(128) = 0$).
\item \textbf{Синхронная инфракрасная термография (Lock-in Thermography):} снижение паразитного тепловыделения решётки более чем на $50\%$ при монохроматической накачке $441\text{ нм}$ по сравнению с референсной ячейкой IBC-$\mathrm{Si}$ ($\hat{n}_E(21) = 0$).
\end{enumerate}

\begin{table}[htbp]
\centering
\caption{Упреждающая диагностика 6 лабораторных режимов срыва и одношаговая матрица коррекции (\texttt{F1}--\texttt{F6}).}
\label{tab:lab-troubleshooting-f1-f6}
\footnotesize
\begin{tabular}{@{}p{0.18\linewidth}p{0.27\linewidth}p{0.24\linewidth}p{0.23\linewidth}@{}}
\toprule
\textbf{Режим срыва} & \textbf{Наблюдаемый симптом} & \textbf{Причина в ДКМ} & \textbf{Одношаговая коррекция} \\
\midrule
\textbf{F1.} Дегидрированный сухой оксид & Люминесценция тетрацена погашена, но $\mathrm{MFE}$ в фототоке $\mathrm{Si}$ равен нулю; $\tau_{\mathrm{eff}} < 200\text{ мкс}$ & Потеря протона $Z : 46 \to 45$ ($45^2 = 2025 < 2048$): пробой затвора Паули синглетом $S_1$ & Отжиг в $\mathrm{H_2}/\mathrm{N_2}$ ($520\text{--}570\text{ К}$) до $[\mathrm{AlO_4H}]/\mathrm{SiO_4}$ ($Z=46,\ 46^2>2048$) \\
\addlinespace
\textbf{F2.} Переутолщённый диэлектрик & Сильный $\mathrm{MFE}$ в люминесценции тетрацена и пик $882\text{ нм}$, но нулевой прирост тока $\mathrm{Si}$ & Мост $>3$ монослоёв ($d_{\mathrm{ox}} > 1{,}5\text{ нм}$) обрывает декстеровский обмен $O_2(3)$ & Удержать толщину моста строго в окне $2\text{--}3$ монослоёв ($0{,}8\text{--}1{,}2\text{ нм}$) \\
\addlinespace
\textbf{F3.} Аморфный напыл или $\mathrm{Si}(111)$ & Сублинейный завал фототока при $\ge 1\text{ Sun}$ (биэкситонное Оже-гашение) & Потеря ориентации $\Delta\varphi = 128$: при $\Delta\varphi \to 0$ перекрытие $\mathcal{C}_3(0) = +128$ включает Оже & Подложка строго $\mathrm{Si}(100)$ ($d_{\mathrm{in}}(\square)=4$), напыление $0{,}2\text{--}0{,}5\text{ \AA/с}$ при $295\text{--}300\text{ К}$ \\
\addlinespace
\textbf{F4.} Переутолщённый слой тетрацена & Зеркальный провал $\mathrm{EQE}$ на $441\text{ нм}$ («анти-спектр поглощения») & При $d_{\mathrm{Tc}} > 100\text{ нм}$ триплеты рождаются дальше диффузионной длины от $\mathrm{Si}$ & Снизить толщину тетрацена до $30\text{--}70\text{ нм}$ с 12-лучевым кварцевым захватом \\
\addlinespace
\textbf{F5.} Фотоокисление на воздухе & Деградация $\mathrm{EQE}$ и $\mathrm{MFE}$ за $10\text{--}20\text{ мин}$ освещения & Образование эндопероксида рушит седло $120 + 8 = 128$ & Инкапсуляция кварцем (Слой~1) в едином вакуумном/$\mathrm{Ar}$ цикле без засветки \\
\addlinespace
\textbf{F6.} Замена тетрацена на пентацен & Нулевая инжекция триплетов в $\mathrm{Si}$ даже при идеальном мосте $Z=46$ & У пентацена $Z=146>128$, триплет $\lambda\approx 1440\text{ нм} > 1107\text{ нм}$ ниже края $\mathrm{Si}$ & Использовать именно тетрацен $\mathrm{C_{18}H_{12}}$ ($Z=120,\ -\det\mathcal{W}_{\mathrm{MS}}=882\text{ нм}$) \\
\bottomrule
\end{tabular}
\end{table}

\section{Интерактивная диагностика с анти-конфабуляционным стеком \texttt{ai-welcome-atlas} + \texttt{ai-atlas-mcp} и сопровождение под NDA}

Полная верифицированная база операторов открыта в пяти репозиториях \texttt{AI-Guiders}:
\begin{itemize}[noitemsep]
\item \url{https://github.com/AI-Guiders/scientific-papers} --- автономные модули Lean~4 (\texttt{RotatorQuadrature.lean}, \texttt{LogicCohomology.lean}, \texttt{PillaiDiophantine.lean}, \texttt{AbcPlanetaryWronskian.lean}, \texttt{ModelTheoreticLimit.lean}, \texttt{CollatzPlanetarySplice.lean});
\item \url{https://github.com/AI-Guiders/discrete-quantum-mechanics} --- ядро ДКМ (\texttt{lean/MtCarrier/UnifiedOperators.lean}) и настоящая серия книг;
\item \url{https://github.com/AI-Guiders/guiders-physics} --- побитово идентичный цифровой двойник ячейки (\texttt{VerletCell72}, \texttt{CarrierCore.cs}) на \texttt{.NET~11};
\item \url{https://github.com/AI-Guiders/ai-welcome-atlas} и \url{https://github.com/AI-Guiders/ai-atlas-mcp} --- база знаний в стандарте Open Knowledge Format и анти-конфабуляционный MCP-сервер (блокирующий галлюцинации языковых моделей и требующий проверки каждого расчёта компилятором Lean~4).
\end{itemize}
При возникновении на конкретной ростовой установке нестандартных отклонений от сигнатур (\texttt{M1})--(\texttt{M4}), связанных со спецификой реактора ALD/PECVD или легирования пластин, экспериментальные группы могут направить автору паспортные параметры ростового цикла для расчёта индивидуальной карты фазово-оболочечной коррекции (включая работу под NDA / CDA / MTA).

% ============================================================================
\chapter{Сводный реестр критических экспериментов и лабораторно-вычислительных гейтов (\texttt{CE-M}, \texttt{CE-F}, \texttt{CE-A}, \texttt{CE-RD})}
\label{ch:lab-critical-experiments-registry}
% ============================================================================

Помимо твердотельного протокола чистой комнаты по бесфононной кремниевой фотовольтаике (Глава~\ref{ch:lab-solar-photovoltaics}), лабораторный контур единого Естествознания включает четыре серии критических фальсифицируемых проверок, исполняемых командами \texttt{python verify\_principles.py --suite all} (\texttt{69/69 PASS}) и \texttt{lake build} (\texttt{38/38 PASS}):
\begin{enumerate}[label=\textbf{Серия~\arabic*.},leftmargin=*]
\item \textbf{Серия \texttt{CE-M} (Устройство носителя, ротатор и квантово-гидродинамический предел):}
  побитовая идентичность ротатора Грея $(\mathcal{C}_3,\mathcal{S}_3)$ на всех $512$ фазах кольца $\mathbb{Z}_{512}$, постоянный разностный вронскиан $\mathcal{W}_3 \equiv 128$, правило Борна из пина переноса \texttt{carry\_out}, отсутствие изолированных магнитных монополей $\operatorname{div}_{\mathrm{cube}} n_f \equiv 0$, постоянная тонкой структуры $\alpha_{\mathrm{fs}}^{-1}(3) = 137{,}035\,999\,187\dots$ ($0{,}47\sigma$ CODATA-2022) и граница Белла--Цирельсона.
\item \textbf{Серия \texttt{CE-F} и \texttt{CE-A} (Симплициальный атлас $\mathcal{C}_{13}$, рождение 86 вихревых плакет из $\mathrm{IV}=263$ и время жизни пара-позитрония):}
  точное воспроизведение кликового спектра $\mathcal{C}_{13}$ ($K_1=13,\ K_2=36,\ K_3=32,\ K_4=8,\ K_{\ge 5}=0$), калибровочной размерности $\dim(\mathrm{U}(1)\times\mathrm{SU}(2)\times\mathrm{SU}(3)) = 1 + 3 + 8 = 12$, космологического отношения гало голономии вакуума $\Omega_{\mathrm{halo}}/\Omega_b = 16/3$ без тёмной материи и времени жизни пара-позитрония $\tau(\text{p-Ps})$.
\item \textbf{Серия \texttt{CE-RD} (Прикладные и реакторные гейты согласования импедансов):}
  когерентный $587$-кратный коллапс кулоновского барьера в 12-лучевом кубооктаэдре $O_2(3)$ при $\Delta\varphi = 139$, квантовый фазовый храповик $\pm 1$ тик на доменной стенке, 4-канальный нильпотентный томограф Неймана $(2I_4 + N_{\mathrm{nil}})^{-1}$ ($N_{\mathrm{nil}}^3 \equiv 0$) и седловой катализ $|A^\ddagger| = 128$.
\end{enumerate}

\input{volumes/include-measurement-processing.tex}
\input{volumes/include-computational-laboratory.tex}
